\originalchapter{中微子混合与振荡}\label{ch:neutrino}
\originalsection{味态与质量态}
弱作用产生 $\nu_e,\nu_\mu,\nu_\tau$ 味态, 传播的定质量本征态记为 $\nu_i$. 取 $|\nu_\alpha\rangle=\sum_iU_{\alpha i}^*|\nu_i\rangle$, 其中 U 为 PMNS 矩阵. 若不同质量态在产生、传播与探测中保持相干, 传播振幅为
\begin{equation}
 \mathcal A_{\alpha\to\beta}(L)=\sum_iU_{\beta i}U_{\alpha i}^*
 \exp[-\ii m_i^2L/(2E)].
\end{equation}
相位来自 $E_i\simeq p+m_i^2/(2p)$ 并取超相对论传播 $t\simeq L$. 整体共有相位不影响概率, 所以振荡只测质量平方差, 不能独自确定绝对质量.

\originalsection{两味公式}
两味混合矩阵用角 $\theta$ 写 $\nu_\alpha=\cos\theta\nu_1+\sin\theta\nu_2$, 正交味为 $-\sin\theta\nu_1+\cos\theta\nu_2$. 振幅相减后,
\begin{equation}\label{eq:oscillation}
 P_{\alpha\to\beta}=\sin^22\theta\sin^2\frac{\Delta m^2L}{4E},
 \quad P_{\alpha\to\alpha}=1-P_{\alpha\to\beta}.
\end{equation}
实用单位中相位为 $1.267\,\Delta m^2[\mathrm{eV}^2]L[\mathrm{km}]/E[\mathrm{GeV}]$. 振荡消失可能来自混合角零、质量简并, 或能量/基线平均把条纹洗掉; 后者不意味着中微子没有质量.

\originalsection{物质效应与实验配对}
电子中微子在普通物质中有额外带电流相干前向散射势 $V=\sqrt2G_Fn_e$. 在两味近似中去掉单位矩阵项后, 有效 Hamiltonian 是真空项加 $\operatorname{diag}(V,0)$. 对角化得
\begin{equation}
 \sin^22\theta_m=\frac{\sin^22\theta}{(\cos2\theta-2EV/\Delta m^2)^2+\sin^22\theta}.
\end{equation}
密度使对角差抵消时混合变大. 反中微子 $V$ 变号, 三味还包括 CP 相位和多个平方差. 太阳、反应堆、加速器与大气中微子的不同 $L/E$ 分别敏感不同尺度, 不应把所有结果硬套同一个两味参数组.

T2K 源题中的离轴几何通过式 \eqref{eq:offaxis} 集中能谱. 对 $L=295\,\mathrm{km}$ 和典型 $|\Delta m^2|\simeq2.4\times10^{-3}\,\mathrm{eV}^2$, 首个真空极大能量约 $0.57\,\mathrm{GeV}$. 离轴角 $2.5^\circ$ 的能谱集中在此量级, 这是选角的物理原因; 并非固定角度对任何束流都产生同样能谱.

\originalsection{Dirac、Majorana 与质量测量}
若引入右手征 sterile 场, 可形成 Dirac 质量, 左右手征组成一个 Dirac 粒子, 可保持总轻子数. Majorana 质量把一个中性场与其电荷共轭结合, 可违反轻子数. 一代 seesaw 质量矩阵
\begin{equation}
 \begin{pmatrix}0&m_D\\m_D&M_R\end{pmatrix}
\end{equation}
在 $M_R\gg m_D$ 时有轻质量 $|m_{\rm light}|\simeq m_D^2/M_R$ 和重质量约 $M_R$. 这是解释小质量的机制候选, 不是振荡实验已经证明的唯一机制.

绝对质量可由 beta 端点的 $m_\beta^2=\sum_i|U_{ei}|^2m_i^2$、宇宙学的质量和约束以及可能的无中微子双 beta 衰变互补研究. 后者在简单轻 Majorana 交换近似中敏感 $|\sum_iU_{ei}^2m_i|$, 相位可产生抵消. 宇宙学结论依模型, 核衰变推断依核矩阵元. 本稿保留课程的概念与方法, 不声称给出截至 2026 年的最新实验数值界限.
